% !TEX root = main.tex
\documentclass[main.tex]{subfiles}
\begin{document}

\section{Distortion Energies and Affine Deformations}
\label{sec:appendix-distortion}

The softened dual $\jointdual_\alpha$ of \S\ref{sec:regularized-joint} attaches each centroid to its own target by a zero-length spring.
In this appendix, we replace this penalty by a general convex \emph{distortion energy} $\distortion$ of the centroids, which yields a family of concave duals containing $\jointdual_\alpha$ as a special case.
We then characterize when such a dual reproduces an affine deformation of the domain exactly, construct the analog of a Dirichlet energy that does so while also handling general deformations of the boundary, and extend this construction to softened volumes.

\subsection{Generalized Softening}
\label{sec:appendix-distortion-dual}

We promote the prescribed centroids of \S\ref{sec:joint-constraints} to unknowns, penalized by a convex energy $\distortion(\tcentroid)$.
The corresponding primal problem reads
\begin{subequations}
  \label{eq:distortion-primal}
  \begin{alignat}{3}
    & \min_{\cell, \tcentroid} \quad & & \sum\nolimits_{i} \int_{\cell_i} \norm{\xx - \qq_i}^2 \, d\xx + \distortion(\tcentroid) \label{eq:distortion-primal-obj} \\
    & \text{s.t.} \quad & \constraintbrace{16pt}{8pt} & \int_{\cell_i} (\xx - \qq_i) \, d\xx = \pvol_i \, (\tcentroid_i - \qq_i) \ \ \forall i , \label{eq:distortion-primal-cons} \\
    & & & \cellvol_i = \pvol_i \ \ \forall i , \label{eq:distortion-primal-vol}
  \end{alignat}
\end{subequations}
where \eqref{eq:distortion-primal-cons} states, together with \eqref{eq:distortion-primal-vol}, that the targets are the centroids of the cells.
Compared to \eqref{eq:joint-lagrangian}, we write the centroid constraint with the prescribed volume $\pvol_i$ in place of the cell volume $\cellvol_i$, which is equivalent under \eqref{eq:distortion-primal-vol}, but moves the targets out of the integrals over the cells.
This choice is necessary for concavity: with the cell volumes, the targets would multiply the partition in the Lagrangian, and the minimization over the targets would no longer separate from the one over the partitions.
We also measure the first moments relative to the sites rather than to the origin, so that the constraint is invariant under a common translation of the domain, the sites and the targets, whatever the cell volumes.
Introducing the multipliers of \S\ref{sec:joint-constraints}, the targets then enter the Lagrangian only through the terms $\smash{2 \sum_i \pvol_i\,\shift_i\tr(\tcentroid_i - \qq_i) + \distortion(\tcentroid)}$, which do not depend on the partition.
The remaining terms coincide with the Lagrangian \eqref{eq:joint-lagrangian} with targets at the sites.
Denoting by $\jointdual(\shift, \weight; \tcentroid)$ the joint dual \eqref{eq:joint-dual} with targets $\tcentroid$, the minimization over the partitions thus yields $\jointdual(\shift, \weight; \qq)$, whose shifted sites \eqref{eq:shifted-site} simplify to $\smash{\dual_i = \qq_i + \shift_i}$ and $\smash{\combweight_i = \weight_i + \norm{\shift_i}^2}$.
The minimization over the targets yields the convex conjugate $\smash{\distortion^\ast(\mathbf{z}) = \sup_{\tcentroid} \sum_i \mathbf{z}_i\tr\tcentroid_i - \distortion(\tcentroid)}$, which leads to the dual objective
\begin{equation}
  \begin{aligned}
    \jointdual_{\distortion}(\shift, \weight)
    &= \min_{\tcentroid} \Bigl[ \jointdual(\shift, \weight; \qq) + 2 \sum\nolimits_i \pvol_i\,\shift_i\tr(\tcentroid_i - \qq_i) + \distortion(\tcentroid) \Bigr] \\
    &= \jointdual(\shift, \weight; \qq) - 2 \sum\nolimits_i \pvol_i\,\shift_i\tr\qq_i - \distortion^\ast\bigl( \{-2\,\pvol_i\,\shift_i\}_i \bigr) ,
  \end{aligned}
  \label{eq:distortion-dual}
\end{equation}
which is concave for any convex $\distortion$, since $\distortion^\ast$ is convex.

For instance, the penalty of \eqref{eq:regularized-primal}, written in terms of the unknown targets, is $\smash{\distortion(\tcentroid) = \alpha \sum_i \pvol_i \norm{\tcentroid_i - \tcentroid^0_i}^2}$ for fixed targets $\smash{\tcentroid^0}$.
With the slacks $\smash{\slack_i = \tcentroid_i - \tcentroid^0_i}$, the constraint \eqref{eq:distortion-primal-cons} reads $\smash{\int_{\cell_i} (\xx - \tcentroid^0_i) \, d\xx = \pvol_i\,\slack_i}$ under \eqref{eq:distortion-primal-vol}, so \eqref{eq:distortion-primal} is exactly the softened primal \eqref{eq:regularized-primal} with targets $\smash{\tcentroid^0}$.
Both problems thus return the same power diagram, with the same generators and shifts, although $\jointdual_\distortion$ and $\jointdual_\alpha$ are different duals of this primal, whose weights differ by $\smash{2\,\shift_i\tr(\tcentroid^0_i - \qq_i)}$.
More generally, any quadratic $\smash{\distortion(\tcentroid) = \tfrac{1}{2} \sum_{i,j} (\tcentroid_i - \boldsymbol{\psi}_i)\tr \mathrm{K}_{ij} (\tcentroid_j - \boldsymbol{\psi}_j)}$, with a positive definite matrix $\mathrm{K}$ and a rest state $\boldsymbol{\psi}$, has the conjugate $\smash{\distortion^\ast(\mathbf{z}) = \sum_i \mathbf{z}_i\tr\boldsymbol{\psi}_i + \tfrac{1}{2} \sum_{i,j} \mathbf{z}_i\tr (\mathrm{K}^{-1})_{ij}\,\mathbf{z}_j}$, so that
\begin{equation}
  \begin{aligned}
    \jointdual_{\distortion}(\shift, \weight) = \jointdual(\shift, \weight; \qq)
    &+ 2 \sum\nolimits_i \pvol_i\,\shift_i\tr(\boldsymbol{\psi}_i - \qq_i) \\
    &- 2 \sum\nolimits_{i,j} \pvol_i\,\shift_i\tr (\mathrm{K}^{-1})_{ij}\,\pvol_j\,\shift_j .
  \end{aligned}
  \label{eq:distortion-quadratic}
\end{equation}
Since $\mathrm{K}^{-1}$ is dense in general, we keep the targets as unknowns in practice and apply Newton's method to the saddle problem in $(\shift, \weight, \tcentroid)$, whose Hessian is sparse whenever $\mathrm{K}$ is.

\paragraph{Optimality conditions.}
By \eqref{eq:joint-gradient} and the envelope theorem, the gradients of \eqref{eq:distortion-dual} are
\begin{equation}
  \partial_{\shift_i} \jointdual_{\distortion} = 2\,\bigl( \pvol_i\,(\tcentroid_i - \qq_i) - \cellvol_i\,(\centroid_i - \qq_i) \bigr) ,
  \qquad
  \partial_{\weight_i} \jointdual_{\distortion} = \pvol_i - \cellvol_i ,
  \label{eq:distortion-gradient}
\end{equation}
where $\tcentroid$ is the minimizer of the bracket in \eqref{eq:distortion-dual}, which is unique for a strictly convex $\distortion$.
Since the targets enter this bracket only through $\smash{2 \sum_i \pvol_i\,\shift_i\tr(\tcentroid_i - \qq_i) + \distortion(\tcentroid)}$, the minimizer verifies, for any multipliers,
\begin{equation}
  \partial_{\tcentroid_i} \distortion(\tcentroid) = -2\,\pvol_i\,\shift_i .
  \label{eq:distortion-target-optimality}
\end{equation}
At a maximizer, the weight gradient in \eqref{eq:distortion-gradient} vanishes, so the volumes are exact, and the shift gradient then vanishes only if the targets coincide with the centroids.
Evaluating \eqref{eq:distortion-target-optimality} at $\smash{\tcentroid = \centroid}$ and solving for the shifts, the generators $\dual = \qq + \shift$ then read
\begin{equation}
  \dual_i = \qq_i - \frac{1}{2\,\pvol_i}\, \partial_{\tcentroid_i} \distortion(\centroid) .
  \label{eq:distortion-generators}
\end{equation}
Each generator is thus its site corrected by the force of $\distortion$ per unit volume, which recovers \eqref{eq:regularized-generators} for the penalty of $\jointdual_\alpha$.

\paragraph{Reduced primal.}
For cells with volumes $\pvol$ and centroids $\centroid$, the transport cost is linear in the centroids, since
\begin{equation}
  \int_{\cell_i} \norm{\xx - \qq_i}^2 d\xx = \int_{\cell_i} \norm{\xx}^2 d\xx - 2\,\pvol_i\,\qq_i\tr\centroid_i + \pvol_i \norm{\qq_i}^2 ,
  \label{eq:distortion-transport}
\end{equation}
where the first terms sum to the constant $\smash{\int_\domain \norm{\xx}^2 d\xx}$ over all cells.
The primal problem \eqref{eq:distortion-primal} thus reduces to
\begin{equation}
  \min_{\centroid \in \achievable} \ \distortion(\centroid) - 2 \sum\nolimits_i \pvol_i\,\qq_i\tr\centroid_i ,
  \label{eq:distortion-reduced}
\end{equation}
which recovers the projection \eqref{eq:regularized-projection} for the penalty of $\jointdual_\alpha$.
The sites $\qq$ only enter through this linear term, which pushes the centroids towards the boundary of $\achievable$, where power diagrams lie (\S\ref{sec:background}), while $\distortion$ pulls them towards its minimizer.
Accordingly, \eqref{eq:distortion-generators} is the outward normal of $\achievable$ at the optimal centroids.
Without the sites, the minimizer may lie in the relative interior of $\achievable$, where no power diagram realizes the centroids.

\subsection{Affine Deformations}
\label{sec:appendix-distortion-affine}

Consider an affine map $\deform(\xx) = \Fmat\,\xx + \boldsymbol{\tau}$ with $\det\Fmat > 0$, a power diagram $\cell$ of $\domain$ with volumes $\pvol$, centroids $\centroid$, and generators $\dual$ with heights $\height$, and the deformed domain $\domain' = \deform(\domain)$.
The cells $\deform(\cell_i)$ tile $\domain'$, with volumes $\smash{\det\Fmat\,\pvol_i}$ and centroids $\smash{\deform(\centroid_i)}$.
Substituting $\smash{\xx = \Fmat^{-1}(\xx' - \boldsymbol{\tau})}$ into the affine functions $\smash{2\,\dual_i\tr\xx - \height_i}$ of \S\ref{sec:background} shows that $\smash{\deform(\cell_i)}$ is the cell of the generator $\smash{\Fmat^{-t}\dual_i}$ with height $\smash{\height_i + 2\,(\Fmat^{-t}\dual_i)\tr\boldsymbol{\tau}}$.
The deformed cells are thus a power diagram, generated by any homothety $\smash{s\,\Fmat^{-t}\dual + \translation}$.
Note that centroids transform by $\Fmat$, as \emph{vectors}, while generators transform by $\smash{\Fmat^{-t}}$, as \emph{covectors}, consistently with their role as normals of $\achievable$.

Next, we maximize $\jointdual_{\distortion'}$ on $\domain'$ with volumes $\smash{\det\Fmat\,\pvol}$, sites $\qq'$, and an energy $\distortion'$ for the deformed domain.
By \eqref{eq:distortion-generators}, and since generators and volumes determine a unique power diagram \cite{Aurenhammer1998Minkowski}, the maximizer is the deformed diagram $\deform(\cell)$ if and only if
\begin{equation}
  \qq'_i - \frac{1}{2 \det\Fmat\,\pvol_i}\, \partial_{\tcentroid_i} \distortion'\bigl( \deform(\centroid) \bigr) = s\,\Fmat^{-t}\dual_i + \translation
  \label{eq:distortion-affine-condition}
\end{equation}
for some $s \!>\! 0$ and $\translation \!\in\! \R^d$.
This condition explains why $\jointdual_\alpha$ fails to reproduce affine deformations, even with sites and targets transported as $\smash{\qq' \!=\! \deform(\qq)}$ and $\smash{\tcentroid' \!=\! \deform(\tcentroid)}$: both terms on the left of \eqref{eq:distortion-affine-condition} then transform by $\Fmat$, so equality with $\smash{\Fmat^{-t}\dual}$ requires $\smash{\Fmat\tr\Fmat \propto \Id}$, \ie a similarity.
Instead, \eqref{eq:distortion-affine-condition} is met by two separate requirements:
\begin{enumerate}
  \item the energy is stationary at the deformed centroids, \ie its gradient vanishes at $\smash{\tcentroid = \deform(\centroid)}$, and
  \item the sites are the deformed generators up to a homothety, \ie $\smash{\qq' = s\,\Fmat^{-t}\dual + \translation}$ with $s \!>\! 0$.
\end{enumerate}
Then the optimal shifts vanish, and the generators of the maximizer are the sites themselves.
For the formulation to apply to general deformations of the boundary, both $\distortion'$ and $\qq'$ must be computed from the deformed boundary alone, and both requirements must hold whenever this deformation happens to be affine.

\subsection{Dirichlet Energy}
\label{sec:appendix-distortion-dirichlet}

We assume that $\domain$ is a polytope, the \emph{cage}, with planar boundary faces $f$ of outward unit normals $\nrm_f$, and that $\domain'$ is given by a continuous map $\deform$ of $\partial\domain$ that is affine on each face.
The reference is a centroidal power diagram of $\domain$ with volumes $\pvol$, centroids $\centroid$ and, choosing the homothety, generators $\dual \!=\! \centroid$, \eg computed by Lloyd iterations.
We denote by $\area_{if}$ the measure of the boundary patch $\smash{\cell_i \cap f}$, by $\foot_{if}$ the orthogonal projection of $\centroid_i$ onto the plane of $f$, and by $\smash{\ell_{if} = \norm{\centroid_i - \foot_{if}}}$ its distance.
On the reference, we define the stiffnesses
\begin{equation}
  \kappa_{ij} = \frac{\area_{ij}}{\norm{\centroid_i - \centroid_j}} ,
  \qquad
  \kappa_{if} = \frac{\area_{if}}{\ell_{if}} ,
  \label{eq:distortion-stiffness}
\end{equation}
where we assume that each centroid lies on the inner side of the planes of its boundary faces, as for a convex cage, so that $\smash{\foot_{if} - \centroid_i = \ell_{if}\,\nrm_f}$.
These stiffnesses are the coefficients of the volume Hessian \eqref{eq:newton-step} up to a factor of two, augmented by one coefficient per boundary patch.

\paragraph{Closure.}
Since the generators are the centroids, the facet $\face_{ij}$ has outward unit normal $\smash{(\centroid_j - \centroid_i) / \norm{\centroid_j - \centroid_i}}$ with respect to $\cell_i$, while the boundary patch $\smash{\cell_i \cap f}$ has outward unit normal $\nrm_f$.
The vector areas of a closed polytope sum to zero, so the boundary of each cell yields
\begin{equation}
  \sum\nolimits_{j \sim i} \kappa_{ij}\,(\centroid_i - \centroid_j) + \sum\nolimits_{f} \kappa_{if}\,(\centroid_i - \foot_{if}) = 0 .
  \label{eq:distortion-closure}
\end{equation}
Each centroid is thus an average of its neighboring centroids and of its feet on the boundary, weighted by \eqref{eq:distortion-stiffness}, \ie the stiffnesses are generalized barycentric coordinates with linear precision.

\paragraph{Energy.}
We define the distortion energy as the discrete Dirichlet energy of the targets with the cage map as boundary data,
\begin{equation}
  \distortion_\deform(\tcentroid) = \frac{\alpha}{2} \sum\nolimits_{\{i,j\}} \kappa_{ij} \norm{\tcentroid_i - \tcentroid_j}^2 + \frac{\alpha}{2} \sum\nolimits_{i,f} \kappa_{if} \norm{\tcentroid_i - \deform(\foot_{if})}^2 ,
  \label{eq:distortion-dirichlet}
\end{equation}
where the first sum runs over pairs of neighbors in the reference, and $\alpha \!>\! 0$ is a stiffness, playing the role of the one in $\jointdual_\alpha$.
This energy is a strictly convex quadratic, since every cell is connected to the boundary through the neighbors of the reference.
Its gradient is
\begin{equation}
  \partial_{\tcentroid_i} \distortion_\deform = \alpha \sum\nolimits_{j \sim i} \kappa_{ij}\,(\tcentroid_i - \tcentroid_j) + \alpha \sum\nolimits_f \kappa_{if}\,\bigl( \tcentroid_i - \deform(\foot_{if}) \bigr) ,
  \label{eq:distortion-dirichlet-gradient}
\end{equation}
so its minimizer $\boldsymbol{\psi}$ solves $d$ sparse linear systems, one per coordinate, and only requires evaluating the cage map at the feet $\foot_{if}$, using the affine extension of $\deform$ from each face $f$ to its plane.
For an affine $\deform$, we have $\smash{\deform(\foot_{if}) - \deform(\centroid_i) = \Fmat\,(\foot_{if} - \centroid_i)}$, so the gradient \eqref{eq:distortion-dirichlet-gradient} at the deformed centroids $\deform(\centroid)$ is $\Fmat$ times the left-hand side of \eqref{eq:distortion-closure}, which vanishes.
The energy \eqref{eq:distortion-dirichlet} therefore meets the first requirement of \S\ref{sec:appendix-distortion-affine}, and for $\deform$ the identity, its minimizer is the reference itself.
For general maps, $\boldsymbol{\psi}$ is a discrete harmonic extension of the cage map evaluated at the centroids, which need not be achievable.
The maximization of $\jointdual_{\distortion_\deform}$ then returns the diagram whose centroids balance this harmonic extension against the push of the sites in \eqref{eq:distortion-reduced}.

\paragraph{Sites.}
The second requirement asks for covectors, which we obtain from the same stiffnesses.
We denote by $\smash{\flux_i = \sum_f \area_{if}\,\nrm_f}$ the vector area of the boundary patches of $\cell_i$, by $\smash{\flux'_i}$ the vector area of their images under $\deform$, and solve
\begin{equation}
  \sum\nolimits_{j \sim i} \kappa_{ij}\,(\qq'_i - \qq'_j) = \flux'_i \quad \forall i
  \label{eq:distortion-sites}
\end{equation}
for the sites $\qq'$.
This is a Neumann problem, solvable since the vector areas of the closed boundary $\partial\domain'$ sum to zero, and its solution is unique up to a translation.
By \eqref{eq:distortion-closure}, the reference centroids solve \eqref{eq:distortion-sites} with $\smash{\flux'_i = \flux_i}$.
For an affine $\deform$, Nanson's formula gives $\smash{\flux'_i = \det\Fmat\,\Fmat^{-t}\flux_i}$, and since \eqref{eq:distortion-sites} acts on each coordinate, the solution is $\smash{\qq' = \det\Fmat\,\Fmat^{-t}\centroid + \translation}$.
These are the deformed generators with dilation $\smash{s = \det\Fmat > 0}$, which meets the second requirement.
Note that $\smash{\flux'_i}$ only depends on the image of the boundary curves of the patches: in 2D, it is the image chord $\smash{\deform(\mathbf{b}) - \deform(\mathbf{a})}$ of each patch from $\mathbf{a}$ to $\mathbf{b}$, rotated clockwise by a quarter turn for a counterclockwise boundary, and in 3D, it is $\smash{\tfrac{1}{2} \oint \deform(\xx) \times d\deform(\xx)}$ along these curves.

\paragraph{Formulation.}
In summary, given the reference and the cage map, we prescribe the volumes $\smash{\pvol_i\,\volume(\domain') / \volume(\domain)}$, compute the sites $\qq'$ from \eqref{eq:distortion-sites}, and maximize $\jointdual_{\distortion_\deform}$ on $\domain'$.
For any affine map and any stiffness $\alpha$, the maximizer is the deformed reference, with zero shifts.
For general maps, its generators \eqref{eq:distortion-generators} are the covector prediction $\qq'$ corrected by the elastic force of the Dirichlet energy per unit volume, while the weights enforce the volumes exactly and act as a pressure.
The stiffness interpolates between the diagram generated by the sites $\qq'$ for $\alpha \!\to\! 0$, which needs no shifts, and the targets $\boldsymbol{\psi}$ for $\alpha \!\to\! \infty$, when they are realizable.
As for $\jointdual_\alpha$, the maximizer is a power diagram only if its generators remain distinct, which holds for small enough $\alpha$ provided the sites $\qq'$ are distinct.
% TODO: quantitative bound on alpha, analog of eq:regularized-distinct.
Finally, the centroidal reference is only needed for the linear precision \eqref{eq:distortion-closure} of the Dirichlet energy.
For a general reference, \eqref{eq:distortion-sites} still applies with $\smash{\kappa_{ij} = \area_{ij} / \norm{\dual_i - \dual_j}}$ and returns the deformed generators, while any energy $\smash{\distortion_\deform(\tcentroid) = \tfrac{1}{2} (\tcentroid - \boldsymbol{\psi})\tr \mathrm{K} (\tcentroid - \boldsymbol{\psi})}$ satisfies the first requirement if its rest state $\boldsymbol{\psi}$ reproduces affine maps, \eg mean value or harmonic coordinates of the reference centroids with respect to the cage.

\subsection{Softened Volumes}
\label{sec:appendix-distortion-volume}

Following Appendix~\ref{sec:softened-volume}, we soften the volume constraints \eqref{eq:distortion-primal-vol} by scalar slacks $\smash{\cellvol_i - \pvol_i = \pvol_i\,\vslack_i}$, penalized by $\smash{\pvol_i\,(\beta\,\vslack_i^2 - \bweight_i\,\vslack_i)}$ in the objective \eqref{eq:distortion-primal-obj}, with stiffness $\beta \!>\! 0$ and reference weights $\bweight$.
The same steps as in Appendix~\ref{sec:softened-volume} yield the concave dual
\begin{equation}
  \jointdual_{\distortion\beta}(\shift, \weight) = \jointdual_{\distortion}(\shift, \weight) - \frac{1}{4\beta} \sum\nolimits_i \pvol_i\,(\weight_i - \bweight_i)^2 ,
  \label{eq:distortion-volume-dual}
\end{equation}
whose optimal weights follow the equation of state \eqref{eq:softened-volume-state}.
The shift gradient in \eqref{eq:distortion-gradient} is unchanged, so its optimal targets are no longer the centroids, but
\begin{equation}
  \tcentroid_i = \qq_i + \frac{\cellvol_i}{\pvol_i}\,(\centroid_i - \qq_i) ,
  \label{eq:distortion-volume-targets}
\end{equation}
which only differ from the centroids by the volume slacks, while the generators still follow \eqref{eq:distortion-generators}, now with the energy evaluated at these targets.
Since the constraint \eqref{eq:distortion-primal-cons} is anchored at the sites, these targets translate with the domain and the sites.
In contrast, anchoring it at the origin would give the targets $\smash{(\cellvol_i / \pvol_i)\,\centroid_i}$, which depend on the choice of origin as soon as the volumes are not exact.

\paragraph{Affine deformations.}
For an affine map $\deform$, the deformed reference has exact volumes and zero shifts, so its targets are its centroids and \eqref{eq:distortion-volume-targets} holds.
Since zero shifts give $\smash{\combweight = \weight}$, the equation of state \eqref{eq:softened-volume-state} with zero volume slacks then requires a third condition in addition to those of \S\ref{sec:appendix-distortion-affine}:
\begin{enumerate}
  \setcounter{enumi}{2}
  \item the reference weights are the weights of the deformed generators $\qq'$, up to an additive constant.
\end{enumerate}
In terms of heights, the reference weights are $\smash{\bweight_i = \norm{\qq'_i}^2 - \height'_i}$, where the heights $\height'$ of the deformed diagram must again be computed from the deformed boundary alone.
We obtain them from a second identity of the stiffnesses \eqref{eq:distortion-stiffness}.
For generators $\dual$ with heights $\height$, the facet $\face_{ij}$ lies on the hyperplane $\smash{2\,(\dual_i - \dual_j)\tr\xx = \height_i - \height_j}$, so its outward unit normal $\smash{\nrm_{ij} = (\dual_j - \dual_i) / \norm{\dual_j - \dual_i}}$ verifies
\begin{equation*}
  \int_{\face_{ij}} \xx\tr\nrm_{ij}\,dA = \frac{1}{2}\,\kappa_{ij}\,(\height_j - \height_i) ,
\end{equation*}
with the stiffness $\smash{\kappa_{ij} = \area_{ij} / \norm{\dual_i - \dual_j}}$ of \eqref{eq:distortion-stiffness}.
The divergence theorem applied to the field $\xx$ over the cell $\cell_i$ then gives
\begin{equation}
  \sum\nolimits_{j \sim i} \kappa_{ij}\,(\height_i - \height_j) = 2\,g_i - 2\,d\,\cellvol_i ,
  \qquad
  g_i = \int_{\partial\cell_i \cap \partial\domain} \xx\tr\nrm \, dA ,
  \label{eq:distortion-heights}
\end{equation}
where $g_i$ equals $d$ times the signed volume of the cones joining the origin to the boundary patches of $\cell_i$.
In 2D, each patch from $\mathbf{a}$ to $\mathbf{b}$ contributes the cross product $\smash{a_1 b_2 - a_2 b_1}$ for a counterclockwise boundary.

We thus compute the heights $\height'$ by solving \eqref{eq:distortion-heights} with the stiffnesses of the reference, the prescribed volumes $\smash{\pvol_i\,\volume(\domain') / \volume(\domain)}$, and the values $g'_i$ of the images of the boundary patches under $\deform$.
This is again a Neumann problem, solvable since its right-hand sides sum to $\smash{2 \int_{\partial\domain'} \xx\tr\nrm\,dA - 2\,d\,\volume(\domain') = 0}$, with a solution unique up to the additive constant allowed by the third condition.
For an affine $\deform$, the sites of \eqref{eq:distortion-sites} are the deformed generators with dilation $\smash{s = \det\Fmat}$.
In this gauge, Nanson's formula scales the facet measures by $\smash{\det\Fmat\,\norm{\Fmat^{-t}\nrm_{ij}}}$ and the generators' distances by the same factor, so the deformed diagram has exactly the stiffnesses of the reference.
Its heights thus verify \eqref{eq:distortion-heights} with the reference stiffnesses, and the solution $\height'$ recovers them up to a constant, which meets the third condition.
Note that the heights do not depend on the translation $\translation$ of the sites, since translating all generators only adds a common linear function to the lifted affine functions, which leaves their heights unchanged.

\paragraph{Formulation.}
In summary, given the reference and the cage map, two sparse Neumann problems with the reference stiffnesses predict the deformed generators $\qq'$ from \eqref{eq:distortion-sites} and their heights $\height'$ from \eqref{eq:distortion-heights}.
We then set the reference weights $\smash{\bweight_i = \norm{\qq'_i}^2 - \height'_i}$ and maximize $\jointdual_{\distortion_\deform\beta}$ on $\domain'$.
For any affine map and any stiffnesses $(\alpha, \beta)$, the maximizer is the deformed reference, with zero shifts and zero volume slacks.
For general maps, $\alpha$ controls the resistance to shape changes through the Dirichlet energy, while $\beta$ controls the compressibility of the cells, with $\beta \!\to\! \infty$ recovering the exact volumes of \S\ref{sec:appendix-distortion-dirichlet}.
For a quadratic energy with a positive definite Hessian, such as \eqref{eq:distortion-dirichlet}, the dual \eqref{eq:distortion-volume-dual} is strongly concave and admits a unique maximizer.

\end{document}
